% Part: lambda-calculus
% Chapter: lambda-definability
% Section: pairs

\documentclass[../../../include/open-logic-section]{subfiles}

\begin{document}

\olfileid{lam}{ldf}{pai}
\olsection{الأزواج ودالة السابق}

\begin{defn}
زوج الحدين~$M$ و$N$، ورمزه~$\tuple{M,N}$، هو بالتعريف
\[
\tuple{M,N} \ident \lambd[f][fMN].
\]
\end{defn}
  
معناه الحدسي دالة تأخذ دالة أخرى فتطبقها في عنصري الزوج.
وعلى هذا نبني حدًا يأخذ حدين ويعطي الزوج المؤلف منهما:
\[
  \fn{Pair} \ident \lambd[mn][\lambd[f][fmn]]
\]
وكما نحتاج إلى تكوين الزوج، نحتاج إلى استخراج عنصريه.
فلذلك نعرّف دالتي وصول، تأخذ كل منهما زوجًا،
فتردّ إحداهما عنصره الأول وتردّ الأخرى عنصره الثاني:
\begin{align*}
  \fn{Fst} & \ident \lambd[p][p(\lambd[mn][m])]\\
  \fn{Snd} & \ident \lambd[p][p(\lambd[mn][n])]
\end{align*}

\begin{prob}
  بيّن وجه أداء دالتي الوصول~$\fn{Fst}$ و$\fn{Snd}$ لما قصد بهما.
\end{prob}

وبعد أن صار تمثيل الأزواج في أيدينا، أمكن أن !!{lambda define} دالة السابق:
\[
  \fn{Pred} \ident \lambd[n][\fn{Fst}(n (\lambd[p][\tuple{\fn{Snd}\, {p}, \fn{Succ}(\fn{Snd}\, {p})}]) \tuple{\num 0, \num 0})]
\]
فقد علمنا أن~$\num n\, f x$ يختزل إلى~$f^{n}(x)$.
أما~$f$ هنا فتأخذ زوجًا~$p$، وتردّ زوجًا مركبته الأولى هي الثانية من~$p$،
ومركبته الثانية خلف تلك المركبة؛ وأما~$x$ فهو~$\tuple{\num 0,\num 0}$.
فتبقى النتيجة~$\tuple{\num 0,\num 0}$ إذا كان~$n=0$،
وتكون~$\tuple{\num{n-1}, \num n}$ فيما عدا ذلك.
ثم تأخذ~$\fn{Pred}$ المركبة الأولى من هذا الناتج.

ومن هنا نعرّف الطرح على الأعداد الطبيعية مع الوقوف عند الصفر؛
فإن جاوز المطروحُ المطروحَ منه فالنتيجة صفر، وإلا فالفرق المعتاد.
وبهذا يتبين قيد العملية التي أطلق الأصل عليها اسم الطرح.
وسبيل تعريفه تكرار تطبيق~$\fn{Pred}$ في~$a$ بمقدار~$b$ من المرات:
\[
\fn{Sub} \ident \lambd[ab][b \fn{Pred}\, a].
\]
    
\end{document}
